\section{Finite length from free amalgamation with a generic order}\label{sec:free-amalg}
\renewcommand{\cal}{\mathcal}

As we saw in Non-example~\ref{non-ex:weak-smooth-approximation}, the ordered atoms do not have oligomorphic approximation. 
Nonetheless they do have the finite length property over any field~\cite[Thm.~4.8]{BFKM24}, not just over those of characteristic zero. 
In this section, we shall generalise that result to a wide class of structures. 
In particular, we will deduce the finite length property of the Rado atoms and its variants.

The structures that we consider here are Fraïssé limits satisfying certain conditions, which are defined by the underlying amalgamation classes.
We shall now state our assumptions and results, with proofs in Section~\ref{sec:cogs-turn}.

\begin{description}
    \item[Graph vocabulary] 
    Consider a finite relational vocabulary $\sigma_0$ consisting of unary and binary relations only. 
    This allows us to talk about, in essence, directed graphs with coloured vertices and edges.

    \item[Free amalgamation class] 
    Let $\Cclass_0$ be a \emph{free} amalgamation class of finite $\sigma_0$-structures. 
    The precise definition may be found in~\cite[Sec.~2.1]{Macpherson11}, 
    but informally it means that when we perform amalgamation, we do not need to glue together any new elements or introduce new relations.
    Here is a useful characterisation~\cite[Lem.~4.5]{SS20} of when an amalgamation class $\Cclass_0$ is free.
    Knowing that $\Cclass_0$ is closed under substructures and isomorphisms, 
    we see that it consists of all the finite $\sigma_0$-structures which do not embed any structure from $\Fclass$, 
    where $\Fclass$ consists of all the minimal (with respect to substructures) finite $\sigma_0$-structures that do not belong to $\Cclass_0$;
    we write \[
        \Cclass_0 = \Forb(\Fclass).
    \]
    That $\Cclass_0$ is a free amalgamation class means that in each $F \in \Fclass$, 
    every two elements $x, y$ are either equal or satisfy at least one of $R(x, y)$ and $R(y, x)$ for some binary relation $R \in \sigma_0$; in that case we will say that $x$ and $y$ are \emph{related}.
    Conversely, and conveniently, given a family $\Fclass$ of finite $\sigma_0$-structures where every two elements are related,
    the class $\Forb(\Fclass)$ of finite $\sigma_0$-structures is a free amalgamation class.    

\end{description}


\begin{example}\label{ex:free-amalg-equality}
    Let $\sigma_0$ be empty. 
    Then $\Cclass_0 = \Forb(\{\})$ is a free amalgamation class consisting of all finite pure sets.
\lipicsEnd\end{example}

\begin{example}\label{ex:free-amalg-graphs}
    Let $\sigma_0$ consist of a single binary relation $E$. 
    A (i)~\emph{graph}, (ii)~\emph{digraph}, (iii)~\emph{tournament} is a $\sigma_0$-structure where $E$ is irreflexive and (i)~symmetric, (ii)~antisymmetric, (iii)~antisymmetric and total.
    For instance, the ordered atoms can be viewed as a tournament once we interpret $E$ as the order relation.

    For more examples, consider $\sigma_0$-structures:
    \[\arraycolsep=20pt\def\arraystretch{1.2}
    \begin{array}{ccccc}
    	\xymatrix@C=15pt{\bigdot\ar@(d,r)[]}
	&
	\quad\xymatrix@C=15pt{\bigdot&\bigdot}
	&
	\xymatrix@C=15pt{\bigdot\ar[r]&\bigdot}
	&
	\xymatrix@C=15pt{\bigdot\ar@<.5ex>[r]&\bigdot\ar@<.5ex>[l]}
	&
	\vcenter{\xymatrix@C=15pt@R=6pt{\bigdot\ar@<.5ex>[dd]\ar@<.5ex>[rd] \\ &\bigdot\ar@<.5ex>[lu]\ar@<.5ex>[ld] \\ \bigdot\ar@<.5ex>[uu]\ar@<.5ex>[ru]}} \\
	\quad A & \quad B & C & K_2 & K_3
	\end{array}
    \]
    %
%    
%    \begin{itemize}
%    \item ${\circlearrowright} = \{0\}$ with $E = \{(0, 0)\}$,
%    \item ${\nrightarrow} = \{1, 2\}$ with $E=\{\}$,
%    \item ${\rightarrow} = \{3, 4\}$ with $E=$
%    \end{itemize}
%    ${\circlearrowright} = \{0\}$, ${\nrightarrow} = \{1, 2\}$, ${\rightarrow} = \{3, 4\}$, ${\rightleftarrows} = \{5, 6\}$, and $\triangle = \{7, 8, 9\}$, 
%    where $E$ is interpreted respectively as $\{(0, 0)\}$, $\{\}$, $\{(3, 4)\}$, $\{(5, 6), (6, 5)\}$, and $\{(7, 8), (8, 7), (8, 9), (9, 8), (9, 7), (7, 9)\}$.
    Then:
    \begin{itemize}
        \item 
        $\Forb(\{A, C\})$ is a free amalgamation class and consists of all finite graphs;
        
        \item
        $\Forb(\{A, C, K_3\})$ is a free amalgamation class and consists of all finite triangle-free graphs;
        
        \item 
        $\Forb(\{A, K_2\})$ is a free amalgamation class and consists of all finite digraphs;

        \item 
        $\Forb(\{A, B, K_2\})$ consists of all finite tournaments. 
        It is an amalgamation class,
        but not a free one: 
        in the minimal forbidden structure $B$, the two elements are not related.
        \lipicsEnd
    \end{itemize}
\end{example}

\begin{description}
    \item[Irreflexivity] 
    For technical reasons, we restrict attention to classes $\Cclass_0$ where all structures are {\em irreflexive}, i.e. such that if $R(x,y)$ then $x\neq y$ for each binary relation $R$.
    % This does not lose generality: for any free amalgamation class $\Cclass_0$, we may replace every binary relation $R$ in the vocabulary with a new unary relation $R_=$, 
    % interpreted in every structure so that $R_=(x)$ if and only if $R(x,x)$, and a binary relation interpreted as the irreflexive part of $R$. 
    % Then $\Cclass_0$ interpreted over this modified vocabulary is still a free amalgamation class, 
    % and all our results transport back to the original $\Cclass_0$. 
    % This way of encoding arbitrary structures as irreflexive ones is standard: see e.g.~\cite[Sec.~2.4]{Herwig1998} and~\cite[p.~121]{SS20}. } 
    This does not lose generality, since we can replace every binary relation $R$ in the vocabulary with its irreflexive fragment $R_{\neq}(x,y)$ and a unary relation $R_=(x)$; see~\cite[Sec.~2.4]{Herwig1998} or~\cite[p.~121]{SS20}.
    All examples considered here are already irreflexive.
    
    \item[Generically ordered expansion] 
    Let $\sigma$ consist of $\sigma_0$ together with a new binary relation $<$.
    Consider the class $\Cclass$ of $\sigma$-structures obtained from $\Cclass_0$ 
    by interpreting $<$ in any $\sigma_0$-structure there as any total order.
    Note that $\Cclass$ is an amalgamation class.
    Indeed, let $X, Y_1, Y_2$ be $\sigma$-structures in $\Cclass$ with $X \subseteq Y_1 \cap Y_2$.
    Then we can amalgamate $Y_1, Y_2$ over $X$ as $\sigma_0$-structures and as $\{<\}$-structures, 
    both using the disjoint union of $Y_1, Y_2$ over $X$ as the underlying set. 
    Superposing these relations gives a $\sigma$-structure in $\Cclass$, which is the desired amalgamation.
    (Because of the total order, this amalgamation in $\Cclass$ is not free except in trivial cases.)
    Denote the Fraïssé limit of $\Cclass$ by $\A$.    
\end{description}




Our main result of this section is that this ordered structure $\A$ has the finite length property over any field.
For clarity, we summarise our assumptions:
\begin{theorem}\label{thm:ordered-free-amalg-has-finite-length}
    Consider:
    \begin{itemize}
        \item $\sigma_0$\quad a finite, at-most-binary relational vocabulary;
        \item $\Cclass_0$\quad a free amalgamation class of finite, irreflexive $\sigma_0$-structures;
        \item $\Cclass$\quad all total orderings of all structures in $\Cclass_0$, over the vocabulary $\sigma_0 \uplus \{<\}$;
        \item $\A$\quad the Fraïssé limit of $\Cclass$~---~a homogeneous and oligomorphic $(\sigma_0 \uplus \{<\})$-structure.
    \end{itemize}
    Then the totally ordered structure $\A$, even with finitely many constants named (i.e.~added as unary relations to the vocabulary), has the finite length property over any field.
\end{theorem}

\begin{remark}
    We do not know how to drop the arity restriction from our proofs, 
    so it is interesting that Evans recently developed a Ramsey-theoretic approach in~\cite{Evans_pm1} which could be used to show that $\Lin_\field \A^2$, 
    for free amalgamation classes over vocabularies of arbitrarily high arity, has finite length.
    We do not know how to combine these results.

    We also do not know whether the finite length property of a general structure $\A$ implies the finite length property of the same structure with finitely many constants fixed.
    There is a quick, positive answer for the equality and the ordered atoms~\cite[Thm.~4.10]{BFKM24},
    but that argument relies on a property that these atoms have~\cite[Lem.~2.21]{BodirskyBodorMarimon_25} and the Rado atoms lack~\cite{no-constant-interpreted-in-Rado}.
\lipicsEnd\end{remark}


 


Before we prove Theorem~\ref{thm:ordered-free-amalg-has-finite-length} in Section~\ref{sec:cogs-turn}, let us state an easy consequence of it and give some examples.
Note that the Fraïssé limit $\A_0$ of $\Cclass_0$ is a first-order reduct of $\A$.
To see this, notice that $\A$, when viewed as a $\sigma_0$-structure, embeds the same finite $\sigma_0$-structures as $\A_0$~---~namely, $\Cclass_0$.
Moreover, it follows from a back-and-forth argument~\cite[Lem.~7.1.4]{hodges1993model} that the $\sigma_0$-structure $\A$ is also homogeneous and therefore isomorphic to $\A_0$.
So we may assume that $\A$ and $\A_0$ have the same underlying set; we then have $\Aut(\A) \subseteq \Aut(\A_0)$, where $\Aut(X)$ denotes the group of automorphisms of $X$.
We thus call $\A$ the \emph{generically ordered expansion} of $\A_0$, or simply the \emph{ordered $\A_0$}.
\begin{corollary}\label{cor:free-finite-length}
    The reduct $\A_0$ of $\A$, even with finitely many constants fixed, has the finite length property over any field.
\end{corollary}
\begin{proof}
    A chain of subspaces in $\Lin_\FF \A_0^d = \Lin_\FF \A^d$ each supported by $S \subset \A_0$ is also a chain of subspaces supported by $S \subset \A$; 
    the latter has a bounded length by the theorem above.
\end{proof}

\begin{example}\label{ex:generically-ordered-equality}
    Continuing from Example~\ref{ex:free-amalg-equality},
    the ordered atoms (Example~\ref{ex:order-atoms}) are the generically ordered expansion of the equality atoms (Example~\ref{ex:equality-atoms}).
\lipicsEnd\end{example}

\begin{example}\label{ex:generically-ordered-graphs}
    Continuing from the first two items of Example~\ref{ex:free-amalg-graphs}, we obtain generically ordered expansions of the Rado graph and of \emph{Henson's triangle-free graph}.
    The ordered Rado graph was studied in~\cite{BPP15} along with its first-order reducts (e.g.~the Fraïssé limit of all finite tournaments). The ordered triangle-free graph will be studied in an extended example in Section~\ref{ex:extended-example}.
\lipicsEnd\end{example}

%So we arrive at the general result below:
\begin{corollary}\label{cor:lachlan-woodrow}
    The following structures have the finite length property over any field:
    \begin{romanenumerate}
        \item the ordered atoms and the equality atoms;
        \item all three \cite{Lachlan1984} countable homogeneous tournaments;
        \item all of the countably many \cite{LachlanWoodrow_80} countable homogeneous graphs, including the Rado graph and Henson's triangle-free graph;
        \item uncountably many \cite[Thm.~2.4]{Henson1972} countable homogeneous digraphs obtained by forbidding tournaments.
    \end{romanenumerate}
\end{corollary}



\section{How the cogs turn: proof of Theorem~\ref{thm:ordered-free-amalg-has-finite-length}}\label{sec:cogs-turn}
To prove Theorem~\ref{thm:ordered-free-amalg-has-finite-length} we will proceed in several steps, 
with the general idea similar to~\cite[Sec.~4.1]{BFKM24}, 
but with significant new complications arising from the presence of non-trivial relations in $\A_0$.
Throughout we fix $\A$ from the statement of Theorem~\ref{thm:ordered-free-amalg-has-finite-length}.

\subsection{Orbits and projections}
To start with, let us view $\A^d$ as $\A^{\{1, \dots, d\}}$. 
More generally, it will be convenient to consider $\A^I$ for a finite totally ordered indexing set $I$ 
(say contained in $\Q$, so that we may take unions later).
Fix a finite support $S \subset \A$.
We shall say that a tuple $a\in\A^I$ is {\em ($S$-)ordered} if $a_i \not\in S$ for all $i$, and $a_i < a_j$ whenever $i < j$. 
Then the orbit $\cal O = \Aut(\A/S) \cdot a$ only contains $S$-ordered tuples, 
and we shall call the orbit $S$-ordered as well.%
\footnote{Here and in the following, $\Aut(\A/S)$ is the (sub)group of those automorphisms of $\A$ that fix every element of $S$. 
So ``$\Aut(\A/S)$-equivariant'' is synonymous with ``supported by $S$''.}
If $a$ is not $S$-ordered, by removing the entries that repeat or come from $S$ and reordering the rest, 
we can always find an $\Aut(\A/S)$-equivariant bijection from $\cal O$ to an $S$-ordered orbit.

To study the lengths of orbit-finitely spanned spaces, we may focus on a single ordered orbit at a time:
\begin{claim}\label{claim:reduction-to-one-orbit}
    The following are equivalent for every finite $S \subset \A$:
    \begin{alphaenumerate}
        \item 
        For each $d \in \{1, 2, \dots\}$, there is some number $l_d$ such that chains of $\Aut(\A/S)$-equivariant subspaces in $\Lin_\FF \A^d$ have length at most $l_d$;
        
        \item 
        For each $d \in \{1, 2, \dots\}$ and every $S$-ordered orbit $\cal O \subseteq \A^d$,
        $\Lin_\FF \cal O$ has finite length with respect to $\Aut(\A/S)$.
    \end{alphaenumerate}
\end{claim}
\begin{claimproof}
    Recall that $\A$ is oligomorphic: given any $S$ and $d$, 
    we know that $\A^d$ is in an $\Aut(\A/S)$-equivariant bijection with a finite disjoint union $\biguplus_i \cal O_i$ of $S$-ordered orbits (possibly in lower dimensions).
    Hence the length of $\Lin_\FF \A^d$, under the action of $\Aut(\A/S)$, equals \[
        \text{length of $\Lin_\FF(\biguplus_i \cal O_i)$} = \text{length of $\bigoplus_i \Lin_\FF \cal O_i$} = \sum_i \text{length of $\Lin_\FF \cal O_i$},
    \]
    which is finite if and only if each summand is finite.
\end{claimproof}

So fix $S$ and an $S$-ordered orbit $\cal O = \Aut(\A/S) \cdot o \subseteq \A^I$.
From here we take an inductive approach.
By $o |^J$ we mean the restriction of $o : I \to \A$ to $J \subseteq I$;
particularly, we will often write $o |^{-i}$ instead of $o |^{ I \setminus \{i\} }$.
The image $\cal O|^J$ of $\cal O$ under this projection agrees with $\Aut(\A/S) \cdot o |^J$ and is still ordered.
The function $(-)|^J$ lifts to a linear $\Aut(\A/S)$-equivariant map
\begin{align*}
    (-)|^J : \Lin_\FF \cal O &\to \Lin_\FF \cal O|^J, 
    \quad \text{defined by} \quad
    v \mapsto \sum_{a \in \cal O|^J} \sum_{b \in \cal O, b|^J = a} v(b) \cdot a.
\end{align*}
Many cancellations can occur under $(-)|^J$. 
Following the terminology used in~\cite[Eq.~(4)]{BFKM24}, 
by the \emph{balanced vectors} of $\Lin_\field \cal O$ we mean the $\Aut(\A/S)$-equivariant subspace below:
\[
    \Ker_\FF \cal O = \bigcap_{i \in I} \ker{ (-)|^{-i} }.
\]

\begin{claim}\label{claim:reduction-to-kernel}
    The following are equivalent for every finite $S \subset \A$:
    \begin{alphaenumerate}
        \item\label{item:LinO-finite-length}
        $\Lin_\FF \cal O$ has finite length with respect to $\Aut(\A/S)$ for every $S$-ordered orbit $\cal O$;

        \item\label{item:balancedLinO-finite-length}
        $\Ker_\FF \cal O$ has finite length with respect to $\Aut(\A/S)$ for every $S$-ordered orbit $\cal O$.
    \end{alphaenumerate}
\end{claim}
\begin{claimproof}
    That (\ref{item:LinO-finite-length}) implies (\ref{item:balancedLinO-finite-length}) is clear, as $\Ker_\FF \cal O \subseteq \Lin_\FF \cal O$.
    
    To prove the other implication, assume (\ref{item:balancedLinO-finite-length}) and let $\cal O \subseteq \A^I$.
    We proceed by induction on $|I|$.
    If $I = \emptyset$, then $\cal O$ must be the entire singleton $\A^\emptyset = \{ () \}$; 
    as $\Lin_\FF \cal O$ has no non-trivial subspaces (let alone finitely supported ones), it has length $1$.
    Now if $|I| \geq 1$, assemble all $|I|$ projection maps into a single map
    \begin{align*}
        \Lin_\FF \cal O \to \bigoplus_{i \in I} \Lin_\FF \cal O|^{-i},
        \quad \text{defined by} \quad
        v \mapsto ( v|^{-i} )_{i \in I},
    \end{align*}
    whose kernel is precisely $\Ker_\FF \cal O$.
    We thus have
    \[
        \text{length of $\Lin_\FF \cal O$}  - \text{length of $\Ker_\FF \cal O$}
        \leq \sum_{i \in I} \text{length of $\Lin_\FF \cal O|^{-i}$}, 
    \]
    which shows that the length of $\Lin_\FF \cal O$ is finite from the assumptions.
\end{claimproof}

As we will see shortly, there exist plenty of balanced vectors.

\subsection{Cogs}
From now on we will use a lightweight notation for combining tuples of atoms: 
for disjoint indexing sets $I$ and $J$ (contained in $\Q$), 
if $a\in\A^I$ and $b\in\A^{J}$ are both ordered, 
then $ab\in \A^{I\cup J}$ will denote their obvious combination. 
We only use this notation if this $ab$ is ordered. 
As an obvious example, for any $S$-ordered $a\in\A^I$ and $J\subseteq I$, we have $a|^{I\setminus J}a|^J=a$.

\begin{definition}\label{def:duo}
    Let $\cal O \subseteq \A^I$ be an $S$-ordered orbit. 
    An \emph{$\cal O$-duo} $a \parallel b$ consists of tuples $a, b \in \cal O$ such that:
    \begin{bracketenumerate}
        \item\label{item:cog-ai-bi} 
        $a_i < b_i$ for all $i\in I$;
        
        \item\label{item:cog-bi-aj}  
        $b_i < a_j$ for all $i<j\in I$;

        \item 
        for any binary $R$ in $\sigma_0$ (which we assumed to be irreflexive) and $i,j\in I$:
        \[
            R(a_i,b_j) \iff R(a_i,a_j) \stackrel{a,b\in\cal O}{\iff} R(b_i,b_j) \iff R(b_i,a_j).           
        \]
     \end{bracketenumerate}
\end{definition}
\begin{remark}\label{rem:duo}
Conditions (\ref{item:cog-ai-bi}) and (\ref{item:cog-bi-aj}) specify a total order on the $2|I|$ atoms in a duo.
Moreover, thanks to irreflexivity, 
each $a_i$ is unrelated to its counterpart $b_i$.
Further, given any $J \subseteq I$, the combined tuple $a|^{I\setminus J} b|^J$ satisfies the same relations as $a$ (and $b$), so it lies in $\cal O$. 
In particular, taking $J=\{i\}$, there is an automorphism $\pi_i$ that sends $a_i$ to $b_i$ and fixes all the other elements of $a$, $b$, and $S$.
\lipicsEnd\end{remark}

For the special case of the ordered atoms (Examples~\ref{ex:order-atoms} and~\ref{ex:generically-ordered-equality}) the following construction was studied in~\cite[p.~11]{BFKM24}, and it already appeared as early as in~\cite[p.~125]{Gray97} under the name of ``polytabloids''.

\begin{definition}
    Given an $\cal O$-duo $a \parallel b$, the corresponding \emph{$\cal O$-cog} is the vector
    \[
        a \between b =
        \sum_{J \subseteq I} (-1)^{|J|}
        (a|^{I \setminus J} b|^{J})
    \]
    in $\Lin_\FF \cal O$. %\footnote{The name ``cog'' turns out to be cognate with Edwards--Venn diagrams.}
    The linear span of all $\cal O$-cogs is denoted by $\Cog_\FF \cal O$.
\end{definition}
Note that, for a fixed $S$-ordered orbit $\cal O$, all $\cal O$-duos (hence all $\cal O$-cogs) are in the same $\Aut(\A/S)$-equivariant orbit.
As a result, $\Cog_\FF \cal O$ is an $\Aut(\A/S)$-equivariant subspace of $\Lin_\FF \cal O$ and it is generated by any single cog.

\begin{claim}\label{claim:cogs-are-balanced}
    $\Cog_\FF \cal O\subseteq \Ker_\FF \cal O$.
\end{claim}
\begin{claimproof}
    Let $\cal O \subseteq \A^I$, let $a\parallel b$ be an $\cal O$-duo, and take any $i \in I$.
    The subsets of $I$ come in pairs of $J$ and $J \cup \{i\}$, where $J\subseteq I \setminus \{i\}$.
    The two tuples $a|^{I\setminus J} b|^{J}$ and $a|^{I \setminus (J \cup \{i\})}b|^{J \cup \{i\}}$ are present in 
    $a\between b$ with the opposite coefficients, and they differ only on the $i$-th entry.
    Therefore they cancel out under $(-)|^{-i}$; hence $(a \between b )|^{-i} = 0$.
\end{claimproof}

It would be remiss not to address the fact that so far, we have not shown that $\cal O$-duos and $\cal O$-cogs even exist in general.
Let us rectify this by showing that they can be found in every non-zero $\Aut(\A/S)$-equivariant subspace of $\Lin_\FF{\cal O}$.


\subsection{Finding cogs}\label{subsec:finding-cogs}
We begin with a technical lemma, which combines the free amalgamation in $\Cclass_0$ and the generic order of $\A$.
\begin{lemma}\label{lem:free-fresh}
    Let $X, Y, \{z\} \subset \A$ be disjoint and finite.
    Then there is an automorphism $\tau \in \Aut(\A)$ such that
    \begin{bracketenumerate}
        \item\label{item:tau-fixes-X} 
        $\tau$ fixes every $x \in X$;

        \item\label{item:tau-frees-from-Yz} 
        $\tau(z)$ is unrelated to all $y \in Y$ and to $z$;
        
        \item\label{item:tau-adds-epsilon} 
        $\tau(z) > z$.
    \end{bracketenumerate}
\end{lemma}
\begin{proof}
Form the free amalgam in ${\Cclass}_0$ of $X\cup Y\cup\{z\}$ and $X\cup\{z\}$ over the common part $X$. The resulting structure can be presented as  $X\cup Y\cup\{z,z'\}$ for some new element $z'$. To make this a $\sigma$-structure, inherit the order on $X \cup Y \cup \{z\}$ from $\A$, and declare that $z < z'$, as well as $z' < a$ whenever $z<a$ for $a\in X \cup Y$. By homogeneity, $X \cup Y \cup \{z, z'\}$ embeds into $\A$ via some $f$ which is the identity on $X \cup Y \cup \{z\}$;
again by homogeneity, we may extend the embedding 
\[          x \in X \mapsto x, \quad z \mapsto f(z')     \] 
to some automorphism $\tau$ that satisfies (\ref{item:tau-fixes-X}), (\ref{item:tau-frees-from-Yz}), and (\ref{item:tau-adds-epsilon}).
\end{proof}
% \begin{proof}
%     Form the free amalgam of structures in ${\Cclass}_0$:
% % https://q.uiver.app/#q=WzAsNCxbMCwxLCJYIl0sWzEsMiwiWCBcXGN1cCBcXHt6XFx9Il0sWzEsMCwiWCBcXGN1cCBZIFxcY3VwIFxce3pcXH0iXSxbMiwxLCJYIFxcY3VwIFkgXFxjdXAgXFx7eiwgeidcXH0iXSxbMCwxLCJcXHN1YnNldGVxIiwxXSxbMCwyLCJcXHN1YnNldGVxIiwxXSxbMSwzLCJ4IFxcaW4gWCBcXG1hcHN0byB4LFxcIHogXFxtYXBzdG8geiciLDEseyJzdHlsZSI6eyJib2R5Ijp7Im5hbWUiOiJkb3R0ZWQifX19XSxbMiwzLCJcXHN1YnNldGVxIiwxLHsic3R5bGUiOnsiYm9keSI6eyJuYW1lIjoiZG90dGVkIn19fV1d
% \[\begin{tikzcd}
% 	& {X \cup Y \cup \{z\}} & \\
% 	X && {X \cup Y \cup \{z, z'\}} \\
% 	& {X \cup \{z\}}
% 	\arrow["\subseteq"{description}, dotted, from=1-2, to=2-3]
% 	\arrow["\subseteq"{description}, from=2-1, to=1-2]
% 	\arrow["\subseteq"{description}, from=2-1, to=3-2]
% 	\arrow["{x \in X \mapsto x,\ z \mapsto z'}"{description}, dotted, from=3-2, to=2-3]
% \end{tikzcd}\]
%     so that no element of $Y \cup \{z\}$ is related to $z'$.
%     To make $X \cup Y \cup \{z, z'\}$ a $\sigma$-structure, 
%     inherit the order on $X \cup Y \cup \{z\}$ from $\A$,
%     and declare that $z < z'$, as well as $z' < a$ whenever $z<a$ for $a\in X \cup Y$. (This makes the above a diagram of embeddings in the presence of the order; we only use the bottom right embedding.)
%     By homogeneity, $X \cup Y \cup \{z, z'\}$ embeds into $\A$ via some $f$ which is the identity on $X \cup Y \cup \{z\}$;
%     again by homogeneity, we may extend the embedding \[ 
%         x \in X \mapsto x, \quad z \mapsto f(z')
%     \] to some automorphism $\tau$ that satisfies (\ref{item:tau-fixes-X}), (\ref{item:tau-frees-from-Yz}), and (\ref{item:tau-adds-epsilon}).
% \end{proof}

\begin{lemma}\label{lem:cog-fresh-single}
    Suppose $a\parallel b$ is an $\cal O$-duo, where $\cal O \subseteq \A^I$ is $S$-ordered.
    Given $z \in S$, 
    \begin{itemize}
        \item write $S' = S \setminus \{z\}$;
        \item let $j \not\in I$ be such that $a z \in \A^{I \cup \{j\}}$~---~thus ${\cal O'}=\Aut(\A / S') \cdot az \subseteq \A^{I \cup \{j\}}$~---~is $S'$-ordered;
        \item let $X \subset \A$ be any finite set containing $\{a_i, b_i \mid i \in I\} \cup S'$ but not $z$.
 %       \item let $Y \subset \A$ be any finite set disjoint from $X \cup \{z\}$;
    \end{itemize}
    Denote $z'=\tau(z)$, where $\tau \in \Aut(\A/X)$ is afforded by Lemma~\ref{lem:free-fresh} (with an arbitrary $Y$). Then $az \parallel bz'$ is an $\cal O'$-duo.
\end{lemma}
\begin{proof}
    First, notice that $bz'\in {\cal O}'$ and that we have the required order relations with $z$ and $z'$. 
    The remaining condition (3) of Definition~\ref{def:duo}, for any binary $R$ in $\sigma_0$, splits into the following cases (and their symmetric counterparts):
\begin{itemize}
\item $R(a_i,b_{i'})\iff R(a_i,a_{i'})$ since $a \parallel b$ is an $\cal O$-duo;
\item $R(a_i,z')\iff R(a_i,z)$ since $\tau$ is an automorphism that fixes all $a_i$;
\item $R(a_i,z)\iff R(b_i,z)$ since $a, b \in \cal O$ and $z\in S$;
\item $R(z,z')$ and $R(z,z)$ are both false: $z'$ is unrelated to $z$ by Lemma~\ref{lem:free-fresh}, and $R$ is irreflexive. \qedhere
\end{itemize}
\end{proof}

Starting from an empty duo, we may apply Lemma~\ref{lem:cog-fresh-single} inductively and obtain:
\begin{lemma}\label{lem:cog-fresh-full}
    Let $\cal O \subseteq \A^I$ be an $S$-ordered orbit.
    Then any $a \in \cal O$ can be extended to an ${\cal O}$-duo $a\parallel b$.
\end{lemma}
% \begin{proof}
%     Enumerate the indices of $I$ as $i_1, \dots, i_d$.
%     Suppose that we have found $b_{i_1}, \dots, b_{i_k}$ such that 
%     \[
%         a|^{\{i_1, \dots, i_k\}} \parallel (i_1 \mapsto b_{i_1}, \dots, i_k \mapsto b_{i_k})
%     \] 
%     is a duo for $\cal O_k = \Aut(\A / S \cup \{a_{i_{k+1}}, \dots, a_{i_d}\}) \cdot a|^{\{i_1, \dots, i_k\}}$~---~note that $() \parallel ()$ is certainly a duo for $\cal O_0$ at the start.
%     If $k < d$, putting $z = a_{i_{k+1}}$, $X = \{a_{i_1}, b_{i_1}, \dots, a_{i_k}, b_{i_k}\} \cup S \cup \{a_{i_{k+2}}, \dots, a_{i_d}\}$, and $Y = \emptyset$, 
%     Lemma~\ref{lem:cog-fresh-single} yields an atom $b_{i_{k+1}}$ that makes
%     \[
%         a|^{\{i_1, \dots, i_k, i_{k+1}\}} \parallel (i_1 \mapsto b_{i_1}, \dots, i_k \mapsto b_{i_k}, i_{k+1} \mapsto b_{i_{k+1}})
%     \] a duo for $\cal O_{k+1}$.
%     For $k=d$ this gives the desired duo for $\cal O_d = \cal O$.
% \end{proof}

As some $a \in \cal O$ always exists, it follows that Claim~\ref{claim:cogs-are-balanced} was not vacuous:
we now know $\Cog_\FF{\cal O}$, and hence $\Ker_\FF{\cal O}$, is non-zero~---~but barely so.
Indeed, as the result below shows, $\Cog_\FF{\cal O}$ admits no non-trivial $\Aut(\A/S)$-equivariant subspaces.
\begin{theorem} \label{thm:cogs-arise-everywhere}
    Any non-zero $\Aut(\A/S)$-equivariant subspace $V\subseteq \Lin_\FF \cal O$ contains $\Cog_\FF \cal O$.
\end{theorem}
\begin{proof}
    Pick any $v \in V$ and $a \in \cal O$ with $v(a)\neq 0$;
    it is enough to show that $V$ contains $a \between b$ for some $b\in \cal O$.
    Define:
    \[
        S^* = S \cup \{c_i \mid v(c) \neq 0, i \in I\} \setminus \{a_i \mid i \in I\} \supseteq S
    \]
    and put $\cal O^* = \Aut(\A / S^*) \cdot a \subseteq \cal O$~---~then $\cal O^*$ is $S^*$-ordered.
    By Lemma~\ref{lem:cog-fresh-full}, we can find $b \in \cal O^*$ such that $a \parallel b$ is an $\cal O^*$-duo and \emph{a fortiori} an $\cal O$-duo.
    Take the automorphisms $\pi_{i_1}, \dots, \pi_{i_d}$ from Remark~\ref{rem:duo}, where $i_1, \dots, i_d$ enumerate $I$.
    Now define $v^{(0)} = v$ and \[
        v^{(k)} = v^{(k-1)} - \pi_{i_k} v^{(k-1)}.
    \]
    Then each $v^{(k)}$ is in $V$.
   By induction on $k$, we have:
    \[
        v^{(k)} = \sum_{c \in C_k} \sum_{J \subseteq \{i_1, \dots, i_k\}} (-1)^{|J|} v(c) \left(\prod_{j \in J} \pi_j c\right),
    \]
    where
    \[C_k = \{ c \mid\ v(c) \neq 0 \text{ and } \{c_{i_1}, \dots, c_{i_k}, \dots, c_{i_d}\} \supseteq \{a_{i_1}, \dots, a_{i_k}\}\ \}.
    \]
    But $\{c_{i_1}, \dots, c_{i_d}\} \supseteq \{a_{i_1}, \dots, a_{i_d}\}$ means that $c = a$, so $C_d=\{a\}$ and $\frac{1}{v(a)}v^{(d)}=a\between b$.
\end{proof}

\begin{corollary}\label{cor:cog-simple}
    $\Cog_\FF \cal O$ has length $1$.
\end{corollary}

In light of Claims~\ref{claim:reduction-to-one-orbit},~\ref{claim:reduction-to-kernel} and~\ref{claim:cogs-are-balanced}, 
for the finite length property it is enough to prove that $\Ker_\FF \cal O \subseteq \Cog_\FF \cal O$. 
In words, we need to show that every balanced vector in $\Lin_\FF{\cal O}$ is a linear combination of $\cal O$-cogs.
Before stating this (as Theorem~\ref{thm:cog-span-generally}), let us illustrate the key ideas of a proof on an example.

\subsection{Spanning by cogs: an extended example}
\label{ex:extended-example}
Let $\A_0$ be the universal triangle-free (undirected) graph, introduced by Henson~\cite{Henson1971}, and let $\A$ be its generically totally ordered version. 
Consider nine atoms $\{a,\ldots,i\}$, ordered by~$<$ alphabetically, with the edge relation as shown here:
\[\scalebox{0.8}{
\rotatebox{7}{
    \xymatrix@R=6pt@C=15pt{
    \rotatebox{-7}{$h$} & & & & & & \rotatebox{-7}{$i$} \\
    \\
    & & & \rotatebox{-7}{$g$}\ar@{-}[uulll]\ar@{-}[uurrr] \\
    & \rotatebox{-7}{$e$} & & & & \rotatebox{-7}{$f$} \\
    & & \rotatebox{-7}{$c$}\ar@{-}[ul]\ar@{-}[uur] & & \rotatebox{-7}{$d$}\ar@{-}[uul]\ar@{-}[ur] \\
    \rotatebox{-7}{$a$}\ar@{-}[rrrrrr]\ar@{-}[uuuuu]\ar@{-}[uur] & & & & & & \rotatebox{-7}{$b$}\ar@{-}[uul]\ar@{-}[uuuuu]
    }
}
}\]
This graph is drawn so that the total order of the atoms corresponds to the vertical order.

Putting $S=\emptyset$ and $|I|=2$, let ${\cal O}$ be the ordered orbit of pairs of atoms which are adjacent. 
Consider the following vector:
\[
    v = ah - ae + ce - cg + dg - df + bf - bi + gi - gh \in \Lin_\FF{\cal O}.
\]
This can be graphically presented as the following graph:
\[\scalebox{0.8}{
\rotatebox{7}{
    \xymatrix@R=6pt@C=15pt{
    \rotatebox{-7}{$h$} & & & & & & \rotatebox{-7}{$i$} \\
    \\
    & & & \rotatebox{-7}{$g$}\ar@[blue][uulll]\ar@[red][uurrr] \\
    & \rotatebox{-7}{$e$} & & & & \rotatebox{-7}{$f$} \\
    & & \rotatebox{-7}{$c$}\ar@[red][ul]\ar@[blue][uur] & & \rotatebox{-7}{$d$}\ar@[red][uul]\ar@[blue][ur] \\
    \rotatebox{-7}{$a$}\ar@{-}[rrrrrr]\ar@[red][uuuuu]\ar@[blue][uur] & & & & & & \rotatebox{-7}{$b$}\ar@[red][uul]\ar@[blue][uuuuu]
    }
}
}\]
where edges with coefficient $+1$ are marked as red, and with $-1$ as blue. 
The arrows on the chosen edges remind us that the pairs in $\cal O$ are ordered, but this is mere decoration: 
the definition of $\cal O$ means that all arrows must point upwards.

Note that $v$ is balanced. 
Graphically, this means that every atom has as many outgoing red edges as outgoing blue edges, and as many incoming red edges as incoming blue edges.

It is easy to draw $\cal O$-cogs in this way. 
Given some additional atom $z>h$ which is adjacent to $a$ and $g$ but not to $h$, the $\cal O$-cog $ah\between gz$ can be drawn as:
\[\rotatebox{15}{
    \xymatrix{
    \rotatebox{-15}{$h$} & \rotatebox{-15}{$z$} \\
    \rotatebox{-15}{$a$}\ar@[red][u]\ar@[blue][ur] & \rotatebox{-15}{$g$}\ar@[blue][ul]\ar@[red][u]  
    }
}\]

We would like to present $v$ as a sum of such ${\cal O}$-cogs. 
Some additional atoms must be used for that, as no four atoms among the original nine form an $\cal O$-duo.
It would be very convenient to introduce a single new atom $z$ to form all the $\cal O$-duos that we will use. 
(Such a new atom is not necessary if $\A$ is the ordered atoms, as then we can simply take $z$ to be the largest atom present in the vector; 
see also the proof of~\cite[\text{Cl.~4.7}]{BFKM24}.)
We can naïvely require $z$ to be:
\begin{itemize}
    \item 
    larger than every atom in $v$,

    \item 
    adjacent to every atom that is a source of a directed edge in $v$ 
    (equivalently: that occurs as the first component of a pair in $v$), and
    
    \item 
    not adjacent to any atom that is a target of a directed edge in $v$ 
    (equivalently: that occurs as the second component of a pair in $v$).
\end{itemize}
However, such a $z$ does not exist in the ordered triangle-free graph $\A$. 
There are two problems:
\begin{itemize}
    \item 
    The atom $g$ occurs both as the first and as the second component in pairs present in $v$. 
    Our specification of whether $z$ is adjacent to $g$ is therefore inconsistent.
    
    \item 
    Atoms $a$ and $b$ both occur as first components in $v$, and they are adjacent in $\A$. 
    As a result, an atom $z$ as prescribed would create a triangle $abz$ in $\A$, which is forbidden.
\end{itemize}
We resolve such {\em conflicts} by considering auxiliary atoms $g'>g$ and $b'>b$, with just enough edges to make $gh\parallel g'i$ and $bf\parallel b'i$ valid ${\cal O}$-duos. 
Specifically, we postulate edges $E(g', h), E(g', i), E(b', f), E(b', i)$ and no more. 
Such atoms exist by Lemma~\ref{lem:free-fresh}. 
We then define:
\[
    v' = v + (gh \between g'i) - (bf \between b'i)
\]
which can be drawn as:
\[\scalebox{0.8}{
\rotatebox{7}{
    \xymatrix@R=6pt@C=15pt{
    \rotatebox{-7}{$h$} & & & & & & \rotatebox{-7}{$i$} \\
    \\
    & & & \rotatebox{-7}{$g$}\ar@{-}[uulll]\ar@{-}[uurrr] & \rotatebox{-7}{$g'$}\ar@[blue]@/_2ex/[uullll]\ar@[red]@/_1ex/[uurr] \\
    & \rotatebox{-7}{$e$} & & & & \rotatebox{-7}{$f$} \\
    & & \rotatebox{-7}{$c$}\ar@[red][ul]\ar@[blue][uur] & & \rotatebox{-7}{$d$}\ar@[red][uul]\ar@[blue][ur] \\
    \rotatebox{-7}{$a$}\ar@{-}[rrrrrr]\ar@[red][uuuuu]\ar@[blue][uur] & & & & & & \rotatebox{-7}{$b$}\ar@{-}[uul]\ar@{-}[uuuuu] & \rotatebox{-7}{$b'$}\ar@[red]@/_1ex/[uull]\ar@[blue]@/_1ex/[uuuuul]
    }
}
}\]
Now an atom $z$ as postulated above does not create any triangles:
\[\scalebox{0.8}{
\rotatebox{7}{
    \xymatrix@R=6pt@C=15pt{
    & & & \rotatebox{-7}{$z$}\ar@{-}@/_4ex/[llldddddd]\ar@{-}[lddddd]\ar@{-}[rddddd]\ar@{-}@/_-0.5ex/[rddd]\ar@{-}@/_-3ex/[rrrrdddddd] \\
    \rotatebox{-7}{$h$} & & & & & & \rotatebox{-7}{$i$} \\
    \\
    & & & \rotatebox{-7}{$g$}\ar@{-}[uulll]\ar@{-}[uurrr] & \rotatebox{-7}{$g'$}\ar@[blue]@/_2ex/[uullll]\ar@[red]@/_1ex/[uurr] \\
    & \rotatebox{-7}{$e$} & & & & \rotatebox{-7}{$f$} \\
    & & \rotatebox{-7}{$c$}\ar@[red][ul]\ar@[blue][uur] & & \rotatebox{-7}{$d$}\ar@[red][uul]\ar@[blue][ur] \\
    \rotatebox{-7}{$a$}\ar@{-}[rrrrrr]\ar@[red][uuuuu]\ar@[blue][uur] & & & & & & \rotatebox{-7}{$b$}\ar@{-}[uul]\ar@{-}[uuuuu] & \rotatebox{-7}{$b'$}\ar@[red]@/_1ex/[uull]\ar@[blue]@/_1ex/[uuuuul]
    }
}
}\]
and it is easy to calculate:
\[
v' = (ah \between g'z)  - (ae \between cz) - (cg \between dz) + (b'f \between dz) -(b'i \between g'z)
\]
which presents $v = v' - (gh \between g'i) + (bf \between b'i)$ as a linear combination of $\cal O$-cogs.



\subsection{All those equivariant subspaces}\label{sec:equivariant-subspaces}
Inspired by the above example, we shall now assert in somewhat greater generality:

\begin{theorem}\label{thm:cog-span-generally}
    Every $v \in \Ker_\EE \cal O$ can be written as
    \[
        v = \sum_{a \parallel b} \lambda_{a \parallel b} \cdot a \between b
    \]
    with $\lambda_{a \parallel b} \in \EE$,
    where $\EE$ is any subspace of some $\FF^n$.
\end{theorem}

\begin{remark}\label{rem:FF-EE}
We must tread carefully here, as $\EE$ is not a field.
For example $\EE$ can be $\{ (\kappa, -\kappa) \mid \kappa \in \FF \}$ in $\FF^2$. 
Then $\Lin_\EE \cal O$ is an equivariant subspace of $\Lin_{\FF^2} \cal O$; 
the latter is more traditionally viewed as $\Lin_\FF \cal O \oplus \Lin_\FF \cal O$.

Given $\lambda \in \EE$ and $a \in \cal O$, we need to understand
\(
    \lambda \cdot a
\)
as a formal expression for an element of $\Lin_\EE \cal O$, rather than the result of a scalar multiplication.
Accordingly, we need to redefine $\Cog_\EE\cal O$ to be spanned by formal expressions $\lambda\cdot(a \between b)$ for $\lambda\in\EE$. 
We still have $\Lin_{\FF^n} \cal O \supseteq \Lin_\EE \cal O \supseteq \Ker_\EE \cal O \supseteq \Cog_\EE \cal O$ as equivariant spaces over $\FF$.
\lipicsEnd\end{remark}

The proof of Theorem~\ref{thm:cog-span-generally} follows the lines of the example in Section~\ref{ex:extended-example}, by carefully removing atom conflicts from a vector before inventing a fresh atom that creates enough duos and cogs. Technical details are deferred to an extended version of this paper.

Putting $\EE = \FF$ in Theorem~\ref{thm:cog-span-generally}, we have~---~as discussed at the end of Section~\ref{subsec:finding-cogs}~---~established Theorem~\ref{thm:ordered-free-amalg-has-finite-length}.
But by allowing $\EE$ to be finite-dimensional vector spaces, we can now describe all equivariant subspaces of an orbit-finite-dimensional vector space: 
we need only look at local sums of coefficients.
%(Stated in a another way, there are no unexpected subspaces here ---
%unlike the infinite chain exhibited in~\cite[Thm.~4.14]{BFKM24} for the bit-vector atoms from Example~\ref{ex:vector-space-atoms}.) 

\begin{theorem}[Compare with~{\cite[Cor.~3.17]{Gray97},~\cite[Thm.~15]{HofmanLerouxTotzke_17}, \cite[Thm.~3.4]{HofmanRozycki_2022}, \cite[Sec.~6]{GhoshHofmanLasota_22}, \cite[Thm.~1.4]{Evans_pm1}}]\label{thm:coeff-approximation}
    Let $\A$ be as in Theorem~\ref{thm:ordered-free-amalg-has-finite-length} and fix $d \in \{1, 2, \dots\}$. 
    Then there exists a finite family of equivariant linear maps of the form
    \[
        \restriction_i : \Lin_\FF \A^d \to \Lin_{\FF^{n_i}} \cal O_i,
    \]
    where $\cal O_i$ is an equivariant ordered orbit of $\A^{d_i}$ with $d_i \leq d$, 
    such that every equivariant subspace $W \subseteq \Lin_\FF \A^d$ is equal to
    \[
        \{v \in \Lin_\FF \A^d \mid\ \forall i, \forall a \in \cal O_i : v{\restriction_i}(a) \in \EE_i\ \}
    \]
    with the finite-dimensional subspaces $\EE_i \subseteq \FF^{n_i}$ given by
    \[
        \EE_i = \{w {\restriction_i}(b) \mid\ w \in W, b \in \cal O_i\ \}.
    \]
\end{theorem}

\begin{example}
    Let $\A$ be the ordered atoms. 
    For $d = 1$ there are two maps, $\restriction_{1} : \Lin_\FF \A \to \Lin_\FF \A$ and $\restriction_{2} : \Lin_\FF \A \to \FF$, given by  
    \[
        v {\restriction_1} (a) = v(a), \qquad
        v {\restriction_2} () = \sum_{x} v(x).
    \]
    For an equivariant vector space $V \subseteq \Lin_\FF \A$, the subspace $V {\restriction_1}(\A) \subseteq \FF$ is either $\{0\}$ or $\FF$.
    \begin{itemize}
        \item In the first case, $V$ must be the zero space.
        \item In the second case, we similarly distinguish two cases for $V {\restriction_2}() \subseteq \FF$:
        \begin{itemize}
            \item if it is the entire $\FF$, then $V$ must be the full space $\Lin_\FF \A$ by Theorem~\ref{thm:coeff-approximation};
            \item if it is $\{0\}$, then $V$ must be the zero-sum space (which exists and is spanned by $\{a - b \mid a, b \in \A\}$) again using Theorem~\ref{thm:coeff-approximation}.
        \end{itemize}
    \end{itemize}
    So $\Lin_\FF \A$ has the same structure as the space over the equality atoms in Example~\ref{ex:length-one-dim-one}.

    For $d = 2$, there are three maps, one of which is of the form $\restriction_2 : \Lin_\FF \A^2 \to \Lin_{\FF^5} \A$. 
    By considering subspaces of $\FF^5$, when the field is infinite we can find infinitely many equivariant subspaces in $\Lin_\FF \A^2$.
    Note that the length of $\Lin_\FF \A^2$ is still finite (and equal to $2^1 + 2^2 + 2^2$, as we will see below).
\lipicsEnd\end{example}

Going through the finite-dimensional subspaces of the $\FF^{n_i}$'s gives a complete list of the equivariant subspaces of $\Lin_\FF \A^d$, by Theorem~\ref{thm:coeff-approximation}, but note that it is far from being irredundant.
In particular, to bound the length of $\Lin_\FF \A^d$ from below, we still need to exhibit a long chain of equivariant subspaces.
It is fortunately straightforward to adapt \cite[Cor.~4.12]{BFKM24} and establish an exact length:
\begin{corollary}\label{cor:2d-length}
    The length of $\Lin_\FF \cal O$, where $\cal O \subseteq \A^d$ is an ordered orbit, is precisely $2^d$.
\end{corollary}

As described in~\cite[Sec.~9]{GhoshLasota_24}, results like Theorem~\ref{thm:coeff-approximation} allow us to decide the solvability of orbit-finite systems of linear equations. 
We briefly repeat the argument.
Checking whether the system $\mathbf{A} \mathbf{x} = \mathbf{b}$ admits a solution amounts to checking whether $\mathbf{b}$ is spanned by the columns of $\mathbf{A}$.
In the orbit-finite setting, we assume that the rows and columns are indexed by $\A^d$ and $\A^{d'}$, that each column has finitely many non-zero entries, and that the definition $j \mapsto \mathbf{A}_{-, j}$ of columns is equivariant.
That is, we ask whether $\mathbf{b} \in \Lin_\FF \A^d$ lies in the span of the $\mathbf{A}_{-, j}$'s. 
By Theorem~\ref{thm:coeff-approximation}, it suffices to check whether $\mathbf{b} {\restriction_i} (\cal O_i) \subseteq \mathbf{A}_{-, j_1} {\restriction_i} (\cal O_i) + \cdots + \mathbf{A}_{-, j_r} {\restriction_i} (\cal O_i)$ for all $i$, in the finite-dimensional spaces $\FF^{n_i}$, having chosen orbit representatives $j_1, \dots, j_r \in \A^{d'}$.